% Generated draft: requires mathematical and source review.
\documentclass[11pt]{article}
\usepackage[margin=25mm]{geometry}
\usepackage{fontspec}
\setmainfont{Latin Modern Roman}
\usepackage{amsmath,amssymb,mathtools,unicode-math}
\setmathfont{Latin Modern Math}
\usepackage{enumitem,needspace,xurl,hyperref}
\setlength{\parindent}{0pt}
\setlength{\parskip}{5pt}
\setlength{\emergencystretch}{3em}
\newcommand{\N}{\mathbb{N}}
\newcommand{\Z}{\mathbb{Z}}
\newcommand{\Q}{\mathbb{Q}}
\newcommand{\R}{\mathbb{R}}
\newcommand{\C}{\mathbb{C}}
\newcommand{\F}{\mathbb{F}}
\newcommand{\A}{\mathbb{A}}
\newcommand{\PP}{\mathbb{P}}
\newcommand{\E}{\mathbb{E}}
\newcommand{\eps}{\varepsilon}
\newcommand{\dd}{\,\mathrm{d}}
\begin{document}
\section*{UnsolvedMath ID: 30003063}
\textbf{Source problem code:} catalogue entry
\subsection{Definitions and mathematical statement}
The source leaves the ground ring and algebra hypotheses unspecified. For a conditional conventional formulation, let \(k\) be a commutative ring and \(A\) an associative unital \(k\)-algebra, and put \(A^e=A\otimes_k A^{\mathrm{op}}\). These conventions require confirmation from the source; in particular, no field, flatness, or projectivity hypothesis over \(k\) is supplied. Work in a setting in which the intended Hochschild cohomology admits the stated identification with derived morphisms.

Let
\[
\cdots\longrightarrow P_2\longrightarrow P_1\longrightarrow P_0\longrightarrow A\longrightarrow0
\]
be an arbitrary projective resolution of the left \(A^e\)-module \(A\). Write \(C_P^n=\operatorname{Hom}_{A^e}(P_n,A)\), with differential induced by the differential of \(P\). Let \(D(A^e)\) denote the derived category of \(A^e\)-modules, let \(\Sigma\) denote cohomological suspension, and set
\[
E^n=\operatorname{Hom}_{D(A^e)}(A,\Sigma^n A).
\]
The intended identifications are
\[
H^n(C_P^\bullet)\cong\operatorname{Ext}_{A^e}^n(A,A)
\cong E^n,
\qquad
\phi_n:\operatorname{HH}^n(A)\xrightarrow{\sim}E^n.
\]
The comparison with the intended definition of \(\operatorname{HH}^n(A)\), including its necessary base hypotheses, remains a source qualification.

The Gerstenhaber bracket is the canonical degree \(-1\) graded Lie bracket on Hochschild cohomology:
\[
[-,-]_G:\operatorname{HH}^p(A)\times\operatorname{HH}^q(A)
\longrightarrow\operatorname{HH}^{p+q-1}(A).
\]
Here the graded Lie convention uses shifted degrees \(p-1\) and \(q-1\). In the ordinary module setting, negative-degree groups are zero, so the bracket of two degree-zero classes is zero.

\textbf{Requested construction.} For every such \(A\) and every such projective resolution \(P\), describe how to compute \([x,y]_G\) using \(P\) and representatives in \(C_P^\bullet\). Also construct intrinsically a family of \(k\)-bilinear maps
\[
\beta_{p,q}:E^p\times E^q\longrightarrow E^{p+q-1}
\qquad(p,q\geq0)
\]
such that, for all \(x\in\operatorname{HH}^p(A)\) and \(y\in\operatorname{HH}^q(A)\),
\[
\beta_{p,q}\bigl(\phi_p(x),\phi_q(y)\bigr)
=\phi_{p+q-1}\bigl([x,y]_G\bigr).
\]
For the zero target in degree \(-1\), this identity means the zero map. The requested intrinsic construction must describe the operation on derived morphisms; merely transporting the bracket through \(\phi\) states the compatibility condition. The source does not specify what categorical structure or auxiliary choices an intrinsic construction may use, nor whether the resolution description must provide an operation on cochains themselves.
\subsection{Short English statement}
Describe the Gerstenhaber bracket using any projective bimodule resolution, and construct a corresponding degree \(-1\) bilinear operation intrinsically on derived morphisms from \(A\) to its shifts. The ground-ring hypotheses and the meaning of ``intrinsically'' require clarification. The supplied record labels the problem partially solved; this draft does not independently verify that assessment.
\subsection{Sources}
\begin{itemize}
\item Local UnsolvedMath registry, supplied complete record for permanent ID \(30003063\), exact source code \texttt{catalogue entry}, title \emph{Intrinsic Gerstenhaber Brackets from Projective Resolutions}. The record belongs to the Oberwolfach Reports Open Problems collection and gives proposed year 2016; no individual proposer is supplied.
\item Supplied report citation: \emph{Hochschild Cohomology in Algebra, Geometry, and Topology} (2016), Oberwolfach Reports, pp.~449--506, DOI \texttt{10.4171/owr/2016/10}. Supplied literature-source label: ``OWR report.'' \url{https://doi.org/10.4171/owr/2016/10}. The collection describes these reports as documenting mathematical workshops and their open-problem discussions.
\item Supplied literature-source label: ``bracket from projective resolutions.'' \url{https://doi.org/10.1016/j.jpaa.2016.10.015}. No authors or article title are supplied in the record.
\item Supplied literature-source label: ``extension interpretation.'' \url{https://doi.org/10.1090/S0002-9949-03-03373-7}. No authors or article title are supplied in the record.
\item Literature assessment dated exactly \texttt{2026-08-21}: supplied dated curation, not independently verified. The record labels the problem ``partially solved'' and reports that the bracket can be computed from an arbitrary projective bimodule resolution for the action of \(\operatorname{HH}^1\), that extension-theoretic descriptions exist, and that the full requested intrinsic bimap is not established in complete generality. The background explicitly calls this a dated literature triage derived from the supplied OWR curation and requests independent verification of time-sensitive status and cited claims.
\item The registry labels the statement \texttt{corrected\_verified}, explaining that mathematical content was preserved while notation, typography, or omitted definitions were repaired for self-contained publication. This is supplied registry metadata, not independent mathematical verification of this draft.
\end{itemize}
\end{document}
